\begin{frame}{잔차 추적을 눈으로: 스템 8개로 직선 근사}
  1차원 데이터: 목표 $y = x$ ($x=1,\dots,10$). 쓸 수 있는 모델은
  \textbf{스텀}(깊이 1: threshold보다 작으면 $a$, 크면 $b$)뿐 —
  스템 하나로는 직선을 절대 못 그린다.
  \texttt{learning\_rate=0.5}로 \textbf{8단계} 추격:
  \vskip0.5em
  {\scriptsize
  \begin{tabular}{@{}lll@{}}
    \toprule
    스템 & threshold / left / right & train MSE \\
    \midrule
    1 & 5 \quad $+3.000$ / $+8.000$ & 11.125 \\
    2 & 3 \quad $+0.500$ / $+3.714$ & 3.826 \\
    3 & 7 \quad $+0.617$ / $+3.143$ & 1.403 \\
    4 & 8 \quad $+0.342$ / $+2.071$ & 0.690 \\
    5 & 2 \quad $-0.729$ / $+0.612$ & 0.385 \\
    6 & 9 \quad $+0.054$ / $+1.230$ & 0.270 \\
    7 & 1 \quad $-0.892$ / $+0.195$ & 0.185 \\
    8 & 1 \quad $-0.446$ / $+0.097$ & 0.163 \\
    \bottomrule
  \end{tabular}}
  \vskip0.6em
  \begin{itemize}
    \item train MSE: 스템 0개(모두 0 예측)의 \textbf{38.5}에서 8단계 만에
      \textbf{0.16}까지
    \item 1단계: 왼쪽 절반 $+3.0\times0.5$, 오른쪽 $+8.0\times0.5$로 상수
      $\to$ 2단계 함수는 \textbf{기울기를 못 그려} 남은 오차가 ``기울기''로
      $\to$ 2단계는 threshold=3을 \textbf{1단계의 왼쪽 구간 내부}에 더해
      그 기울기를 수정
    \item tree 7, 8은 이전 트리가 이미 쓴 threshold(=1)를 \textbf{다시} 골라
      진폭만 미세 조정($-0.892\to-0.446$) — ``새 트리가 항상 새 위치''가
      아니라 \textbf{``남은 잔차를 가장 잘 분기하는 곳''}
  \end{itemize}
\end{frame}
